\documentclass[10pt,a4paper]{amsart}
\usepackage[margin=19mm]{geometry}
\usepackage{amsmath,amssymb,mathtools}
\usepackage[scr=rsfs,cal=euler]{mathalfa}
\usepackage{xfrac}
\usepackage[unicode,hidelinks,bookmarks=false]{hyperref}
\newcommand{\ie}{i.e.\ }
\newcommand{\cf}{cf.\ }
\newcommand{\resp}{respectively\ }
\newcommand{\Subsection}{\subsection}
\providecommand{\nobreakdash}{\nobreak\hbox{-}}
\setlength{\emergencystretch}{2em}
\multlinegap0pt
\usepackage{bm}
\usepackage{bbm}
\usepackage{dsfont}
\usepackage{xspace}
\usepackage{cancel}
\usepackage{mathtools}
\usepackage[shortlabels]{enumitem}
\usepackage{amsthm}
\makeatletter
\@ifundefined{theorem}{\newtheorem{theorem}{Theorem}}{}
\@ifundefined{lemma}{\newtheorem{lemma}[theorem]{Lemma}}{}
\@ifundefined{proposition}{\newtheorem{proposition}[theorem]{Proposition}}{}
\makeatother
\title[On the deep commuting graph of a finite group]{On the deep commuting graph of a finite group}
\author{Sumana Hatui, Sanjay Mukherjee, Kamal Lochan Patra}
\date{Source: arXiv:2511.13303v1 (CC BY 4.0)}
\begin{document}
\maketitle

\noindent\emph{Source and licence.} On the deep commuting graph of a finite group. Version arXiv:2511.13303v1 (CC BY 4.0), distributed under the Creative Commons Attribution 4.0 International licence, as recorded in the arXiv licence metadata for this exact version and shown on the abstract page https://arxiv.org/abs/2511.13303v1. Full rights notice: licenses/arxiv-2511.13303-CC-BY-4.0.txt.\par\smallskip

\noindent\textit{Editorial scope.} Counted unit: the Proposition whose source label is \texttt{deep commuting of A\_6,7} in the article (source0.tex lines 786--788, source label \texttt{deep commuting of A\_6,7}), delivered together with the article's own complete original proof of it (source0.tex lines 790--842, one \texttt{proof} environment of 53 lines). No mathematics is written, repaired or re-derived: statement and proof are the article's own text line for line. Required context retained uncounted: disjoint\_commutators (lines 669--675); induced\_S\_n and A\_n (lines 683--690); deep commuting of S\_4,5 (lines 718--720); 3-cycle A\_7 (lines 769--771); 2-cycle A\_7 (lines 778--780). Every definition, display and label of those lines is retained unchanged. Omitted from this package and not redistributed: 1--668, 676--682, 691--717, 721--768, 772--777, 781--785, 789--789, 843--1067. Nothing inside a retained excerpt is omitted or altered: every retained body is delivered exactly as the article's lines are written, so no omission receipt of any kind accompanies this package. Citations: the retained excerpts carry no citation command at all, so no bibliography entry of the article is transcribed and no reference is redistributed. Reference adaptation: none was needed and none was made; every reference inside the delivered unit resolves inside it, so no unresolved reference or citation is produced. Printed slips kept literal: no printed word, symbol, hypothesis or number of the counted statement or of its proof was altered. Layout and class: the article's own class is \texttt{article} with packages including setspace, amsfonts, epsfig, latexsym, amsthm, amsmath, stmaryrd, amssymb, array, enumerate, tikz, tikz-cd, graphicx, graphics; the wrapper is the lane's pinned \texttt{amsart} editorial wrapper at 10pt on A4, which restates the theorem-like environments the retained text needs and adds the article's own macro definitions that the retained text needs (none), with extra wrapper packages (bm, bbm, dsfont, xspace, cancel, mathtools, enumitem). The delivered numbering is the unit's own. Assets: no retained span contains a figure environment, a graphics-inclusion command, a PDF-page inclusion, a table or a TikZ picture; no image file is copied into this package, the package declares no asset dependency and no image file is redistributed with it. Licence: version arXiv:2511.13303v1 of this article is distributed under the Creative Commons Attribution 4.0 International licence, as recorded in the arXiv licence metadata for this exact version and shown on the abstract page https://arxiv.org/abs/2511.13303v1. A full rights notice is delivered at licenses/arxiv-2511.13303-CC-BY-4.0.txt. Characters: every retained excerpt is pure ASCII, so the delivered engine, XeLaTeX with its default Latin Modern text font, reports no missing character.
\begin{lemma}\label{disjoint_commutators}
Let $s_1,s_2,s_3\in S_n$ (or $A_n$) with $\tilde{s_1},\tilde{s_2},\tilde{s_3}$ being their preimages in $\tilde{S_n}$ (or $\tilde{A_n}$), respectively. Then the following happens:
\begin{enumerate}[\rm (i)]
\item If $[s_1,s_3]=e$, then $[\tilde{s_1}\tilde{s_2},\tilde{s_3}]=[\tilde{s_1},\tilde{s_3}][\tilde{s_2},\tilde{s_3}]$;
\item If $[s_1,s_2]=e$, then $[\tilde{s_1},\tilde{s_2}\tilde{s_3}]=[\tilde{s_1},\tilde{s_2}][\tilde{s_1},\tilde{s_3}]$.
\end{enumerate}    
\end{lemma}

\begin{proposition}\label{induced_S_n and A_n}
 For $4\leq m< n$, consider the subgroup $S_m$ of $S_n$ being the group of permutations of the first $m$ symbols inside $S_n$ and the subgroup $A_m$ of $A_n$ being the group of even permutations of the first $m$ symbols inside $A_n$. Then the following happens:
 \begin{enumerate}[\rm (i)]
     \item $\Delta_D(S_n)\vert_{S_m}=\Delta_D(S_m)$;
     \item For $m,n\notin\{6,7\}$, we have $\Delta_D(A_n)\vert_{A_m}=\Delta_D(A_m)$;
     \item $\Delta_D(A_7)\vert_{A_6}=\Delta_D(A_6)$.
 \end{enumerate}    
\end{proposition}

\begin{proposition}\label{deep commuting of S_4,5}
For $n\in\{4,5\}$, we have $\Delta_D(S_n)=\mathcal{P}_e(S_n)$.
\end{proposition}

\begin{lemma}\label{3-cycle A_7}
Let $\sigma_1$ and $\sigma_2$ be two disjoint $3$-cycles in $A_7$. Then $\sigma_1\nsim\sigma_2$ in $\Delta_D(A_7)$.    
\end{lemma}

\begin{lemma}\label{2-cycle A_7}
Let $s_1$ and $s_2$ be two distinct permutations of order $2$ in $A_7$ such that $[s_1,s_2]=e$. Then $s_1\nsim s_2$ in $\Delta_D(A_7)$.    
\end{lemma}

\begin{proposition}\label{deep commuting of A_6,7}
For $n\in\{6,7\}$, we have $\Delta_D(A_n)=\mathcal{P}_e(A_n)$.    
\end{proposition}

\begin{proof}
The idea of this proof is similar to the proof of Proposition \ref{deep commuting of S_4,5}. We claim that for any permutation $\sigma\in A_7$, $\mathcal{N}_{\Delta_D(A_7)}(\sigma)=\mathcal{N}_{\mathcal{P}_e(A_7)}(\sigma)$, which yields $\Delta_D(A_7)=\mathcal{P}_e(A_7)$. To show the above claim, we consider the following cases:  

\vspace{.3cm}
\noindent\underline{\bf{Case (i) ($\sigma=(a,b,c)$):}} Consider the permutation $(123)$. Then \[C_{A_7}((123))=\langle (123)\rangle\times\langle (456),(45)(67)\rangle.\] Clearly, $\{e,(132),(45)(67),(46)(57),(47)(56),(123)(45)(67),(123)(46)(57),(123)(47)(56),\\ (132)(45)(67),(132)(46)(57),(132)(47)(56)\}=\mathcal{N}_{\mathcal{P}_e(A_7)}((123))\subseteq \mathcal{N}_{\Delta_D(A_7)}((123))$. Then, any permutation in $C_{A_7}((123))\setminus\mathcal{N}_{\mathcal{P}_e(A_7)}[(123)]$ is one of the form $\lambda, (123)\lambda$ and $(132)\lambda$, where $\lambda$ is a $3$-cycle disjoint from $(123)$. Now from Lemma \ref{disjoint_commutators}, we have $[\tilde{(123)},\tilde{\lambda}]=[\tilde{(123)},\tilde{(123)}\tilde{\lambda}]=[\tilde{(123)},\tilde{(132)}\tilde{\lambda}]$, and from Lemma \ref{3-cycle A_7}
we have $[\tilde{(123)},\tilde{\lambda}]\neq e$. Hence, $\mathcal{N}_{\Delta_D(A_7)}((123))=\mathcal{N}_{\mathcal{P}_e(A_7)}((123))$. Similarly, it can be shown that for any arbitrary permutation of the form $(a,b,c)\in A_7$, we have $\mathcal{N}_{\Delta_D(A_7)}((a,b,c))=\mathcal{N}_{\mathcal{P}_e(A_7)}((a,b,c))$.

\vspace{.3cm}
\noindent\underline{\bf{Case (ii) ($\sigma=(a,b)(c,d)$):}} Consider the permutation $(12)(34)$. Then
\begin{center}
$C_{A_7}((12)(34))=$\\
$(\langle (12)(34),(13)(24)\rangle\times\langle (567)\rangle)\cup(\{(12),(34),(1324),(1423)\}\times\{(56),(57),(67)\}).$   
\end{center}
\vspace{.15cm}
Clearly, $\langle(12)(34)(567)\rangle\cup(\{(1324),(1423)\}\times\{(56),(57),(67)\})=\mathcal{N}_{\mathcal{P}_e(A_7)}[(12)(34)]\subseteq \\ \mathcal{N}_{\Delta_D(A_7)}[(12)(34)]$. Hence, $C_{A_7}((12)(34))\setminus\mathcal{N}_{\mathcal{P}_e(A_7)}[(12)(34)]=(\{(13)(24),(14)(23)\}\times\langle(567)\rangle)\cup(\{(12),(34)\}\times\{(56),(57),(67)\})$. Now, since $(567),(576)\in\mathcal{N}_{\Delta_D(A_7)}[(12)(34)]$, we have $[\tilde{(12)(34)},\tilde{(567)}]=[\tilde{(12)(34)},\tilde{(576)}]=e$. Then from Lemma \ref{disjoint_commutators}, we can say that for any $v\in (\{(13)(24),(14)(23)\}\times\langle(567)\rangle)\cup(\{(12),(34)\}\times\{(56),(57),(67)\})$, $[\tilde{(12)(34)},\tilde{v}]=[\tilde{(12)(34)},\tilde{\xi}]$, for some $\xi\in\{(13)(24),(14)(23)\}\cup(\{(12),(34)\}\times\{(56),(57),(67)\})$. Now, from Lemma \ref{2-cycle A_7} we have $[\tilde{(12)(34)},\tilde{\xi}]\neq e$. Hence, $\mathcal{N}_{\Delta_D(A_7)}((12)(34))=\mathcal{N}_{\mathcal{P}_e(A_7)}((12)(34))$. Similarly, it follows that for any arbitrary permutation of the form $(a,b)(c,d)\in A_7$, we have $\mathcal{N}_{\Delta_D(A_7)}((a,b)(c,d))=\mathcal{N}_{\mathcal{P}_e(A_7)}((a,b)(c,d))$.  

\vspace{.3cm}
\noindent\underline{\bf{Case (iii) ($\sigma=(a,b,c)(d,e,f)$):}} Consider the permutation $(123)(456)$. Then \[C_{A_7}((123)(456))=\langle(123)\rangle\times\langle(456)\rangle.\] Here, $\langle(123)(456)\rangle=\mathcal{N}_{\mathcal{P}_e(A_7)}[(123)(456)]\subseteq\mathcal{N}_{\Delta_D(A_7)}[(123)(456)]$. For the rest of the permutations in $C_{A_7}((123)(456))$, using Lemma \ref{disjoint_commutators} and \ref{3-cycle A_7} we can say the following.
\begin{itemize}
    \item $[\tilde{(123)}\tilde{(456)},\tilde{(123)}]=[\tilde{(456)},\tilde{(123)}]\neq e.$
    \item $[\tilde{(123)}\tilde{(456)},\tilde{(132)}]=[\tilde{(456)},\tilde{(132)}]\neq e.$
    \item $[\tilde{(123)}\tilde{(456)},\tilde{(456)}]=[\tilde{(123)},\tilde{(456)}]\neq e.$
    \item $[\tilde{(123)}\tilde{(456)},\tilde{(465)}]=[\tilde{(123)},\tilde{(465)}]\neq e.$
    \item $[\tilde{(123)}\tilde{(456)},\tilde{(123)}\tilde{(465)}]=[\tilde{(123)},\tilde{(465)}][\tilde{(456)},\tilde{(123)}]=[\tilde{(123)},\tilde{(456)}]\neq e.$
    \item $[\tilde{(123)}\tilde{(456)},\tilde{(132)}\tilde{(456)}]=[\tilde{(123)},\tilde{(456)}][\tilde{(456)},\tilde{(132)}]=[\tilde{(456)},\tilde{(123)}]\neq e.$
\end{itemize}
Hence, we have $\mathcal{N}_{\Delta_D(A_7)}((123)(456))=\mathcal{N}_{\mathcal{P}_e(A_7)}((123)(456))$.

Similarly, it can be shown that for any arbitrary permutation of the form $(a,b,c)(d,e,f)\in A_7$, we have $\mathcal{N}_{\Delta_D(A_7)}((a,b,c)(d,e,f))=\mathcal{N}_{\mathcal{P}_e(A_7)}((a,b,c)(d,e,f))$.

\vspace{.3cm}
\noindent\underline{\bf{Case (iv) ($\sigma=(a,b,c)(d,e)(f,g)$):}} Consider the permutation $(123)(45)(67)$. Then \[C_{A_7}((123)(45)(67))=\langle(123)\rangle\times\langle(45)(67),(46)(57)\rangle.\] Here, $\langle(123)(45)(67)\rangle=\mathcal{N}_{\mathcal{P}_e(A_7)}[(123)(45)(67)]\subseteq\mathcal{N}_{\Delta_D(A_7)}[(123)(45)(67)]$. Now, for any $u\in \langle(123),(132)\rangle$ and $v\in \langle(46)(57),(47)(56)\rangle$, we have $v\in\mathcal{N}_{\mathcal{P}_e(A_7)}[u]\subseteq\mathcal{N}_{\Delta_D(A_7)}[u]$. Hence, $[\tilde{u},\tilde{v}]=e$. Now, from Lemma \ref{disjoint_commutators} we can see that 
\begin{itemize}
    \item $[\tilde{(123)}\tilde{(45)(67)},\tilde{(46)(57)}]=[\tilde{(45)(67)},\tilde{(46)(57)}].$
    \item $[\tilde{(123)}\tilde{(45)(67)},\tilde{(123)}\tilde{(46)(57)}]=[\tilde{(45)(67)},\tilde{(46)(57)}]$.
    \item $[\tilde{(123)}\tilde{(45)(67)},\tilde{(132)}\tilde{(46)(57)}]=[\tilde{(45)(67)},\tilde{(46)(57)}].$
    \item $[\tilde{(123)}\tilde{(45)(67)},\tilde{(47)(56)}]=[\tilde{(45)(67)},\tilde{(47)(56)}]$.
    \item $[\tilde{(123)}\tilde{(45)(67)},\tilde{(123)}\tilde{(47)(56)}]=[\tilde{(45)(67)},\tilde{(47)(56)}]$.
    \item $[\tilde{(123)}\tilde{(45)(67)},\tilde{(132)}\tilde{(47)(56)}]=[\tilde{(45)(67)},\tilde{(47)(56)}]$.
\end{itemize}
In addition, Lemma \ref{2-cycle A_7} yields $[\tilde{(45)(67)},\tilde{(46)(57)}]\neq e$ and $[\tilde{(45)(67)},\tilde{(47)(56)}]\neq e$. Hence, we have $\mathcal{N}_{\Delta_D(A_7)}((123)(45)(67))=\mathcal{N}_{\mathcal{P}_e(A_7)}((123)(45)(67))$. Now, the rest follows via similar arguments.

\vspace{.3cm}
\noindent\underline{\bf{Case (v) ($\sigma=(a,b,c,d)(e,f)$):}} Consider the permutation $(1234)(56)$. Then \[C_{A_7}((1234)(56))=\langle(1234)(56)\rangle.\] Here, $\mathcal{N}_{\mathcal{P}_e(A_7)}[(1234)(56)]=C_{A_7}((1234)(56))$. Hence, it follows that $\mathcal{N}_{\mathcal{P}_e(A_7)}((1234)(56))=\mathcal{N}_{\Delta_D(A_7)}((1234)(56))$. Now, the rest follows similarly.

\vspace{.3cm}
\noindent\underline{\bf{Case (vi) ($\sigma=(a,b,c,d,e)$):}} Consider the permutation $(12345)$. Then \[C_{A_7}((12345))=\langle(12345)\rangle.\] Here, $\mathcal{N}_{\mathcal{P}_e(A_7)}[(12345)]=C_{A_7}((12345))$. Hence, it is not difficult to see that $\mathcal{N}_{\mathcal{P}_e(A_7)}((12345))=\mathcal{N}_{\Delta_D(A_7)}((12345))$. The rest follows similarly.

\vspace{.3cm}
\noindent\underline{\bf{Case (vii) ($\sigma=(a,b,c,d,e,f,g)$):}} Consider the permutation $(1234567)$. Then \[C_{A_7}((1234567))=\langle(1234567)\rangle.\] Here, $\mathcal{N}_{\mathcal{P}_e(A_7)}[(1234567)]=C_{A_7}((1234567))$. Hence, it follows that $\mathcal{N}_{\mathcal{P}_e(A_7)}((1234567))=\mathcal{N}_{\Delta_D(A_7)}((1234567))$. Again, the rest follows similarly as above.

Hence, we have $\Delta_D(A_7)=\mathcal{P}_e(A_7)$. Now, from Proposition \ref{induced_S_n and A_n}(iii) and the fact that the induced subgraph of $\mathcal{P}_e(G)\vert_S$, for any subgroup $S$ of $G$ coincides with $\mathcal{P}_e(S)$, the equality can be deduced for $A_6$. This concludes the proof. 
\end{proof}

\end{document}
